\documentclass[10pt,a4paper]{amsart}
\usepackage{amsmath,amssymb,amsthm,graphicx}
% Retained minimal preamble for the Stacks Project chapter sites-cohomology.tex (Cohomology on Sites).
% Verbatim excerpt of the chapter preamble: Xy-pic setup (native preamble lines 14-18) and the native
% theorem declarations with the mathematical macros (native preamble lines 154-217). Copied mechanically;
% nothing added, removed or redefined. The wrapper class and the amsmath/amssymb/amsthm/graphicx base
% packages are supplied by the assembler, so the class/externaldocument part of the native preamble is omitted.
% For commutative diagrams we use Xy-pic
\usepackage[all]{xy}

% We use 2cell for 2-commutative diagrams.
\xyoption{2cell}
% Theorem environments.
%
\theoremstyle{plain}
\newtheorem{theorem}[subsection]{Theorem}
\newtheorem{proposition}[subsection]{Proposition}
\newtheorem{lemma}[subsection]{Lemma}

\theoremstyle{definition}
\newtheorem{definition}[subsection]{Definition}
\newtheorem{example}[subsection]{Example}
\newtheorem{exercise}[subsection]{Exercise}
\newtheorem{situation}[subsection]{Situation}

\theoremstyle{remark}
\newtheorem{remark}[subsection]{Remark}
\newtheorem{remarks}[subsection]{Remarks}

\numberwithin{equation}{subsection}

% Macros
%
\def\lim{\mathop{\mathrm{lim}}\nolimits}
\def\colim{\mathop{\mathrm{colim}}\nolimits}
\def\Spec{\mathop{\mathrm{Spec}}}
\def\Hom{\mathop{\mathrm{Hom}}\nolimits}
\def\Ext{\mathop{\mathrm{Ext}}\nolimits}
\def\SheafHom{\mathop{\mathcal{H}\!\mathit{om}}\nolimits}
\def\SheafExt{\mathop{\mathcal{E}\!\mathit{xt}}\nolimits}
\def\Sch{\mathit{Sch}}
\def\Mor{\mathop{\mathrm{Mor}}\nolimits}
\def\Ob{\mathop{\mathrm{Ob}}\nolimits}
\def\Sh{\mathop{\mathit{Sh}}\nolimits}
\def\NL{\mathop{N\!L}\nolimits}
\def\CH{\mathop{\mathrm{CH}}\nolimits}
\def\proetale{{pro\text{-}\acute{e}tale}}
\def\etale{{\acute{e}tale}}
\def\QCoh{\mathit{QCoh}}
\def\Ker{\mathop{\mathrm{Ker}}}
\def\Im{\mathop{\mathrm{Im}}}
\def\Coker{\mathop{\mathrm{Coker}}}
\def\Coim{\mathop{\mathrm{Coim}}}

% Boxtimes
%
\DeclareMathSymbol{\boxtimes}{\mathbin}{AMSa}{"02}

%
% Macros for moduli stacks/spaces
%
\def\QCohstack{\mathcal{QC}\!\mathit{oh}}
\def\Cohstack{\mathcal{C}\!\mathit{oh}}
\def\Spacesstack{\mathcal{S}\!\mathit{paces}}
\def\Quotfunctor{\mathrm{Quot}}
\def\Hilbfunctor{\mathrm{Hilb}}
\def\Curvesstack{\mathcal{C}\!\mathit{urves}}
\def\Polarizedstack{\mathcal{P}\!\mathit{olarized}}
\def\Complexesstack{\mathcal{C}\!\mathit{omplexes}}
% \Pic is the operator that assigns to X its picard group, usage \Pic(X)
% \Picardstack_{X/B} denotes the Picard stack of X over B
% \Picardfunctor_{X/B} denotes the Picard functor of X over B
\def\Pic{\mathop{\mathrm{Pic}}\nolimits}
\def\Picardstack{\mathcal{P}\!\mathit{ic}}
\def\Picardfunctor{\mathrm{Pic}}
\def\Deformationcategory{\mathcal{D}\!\mathit{ef}}

% --- appended 2026-10-12 (Chief): the project own definitions of the reference/slogan/history
% display environments, copied mechanically from the shared Stacks preamble lines 8-12
% (\usepackage{verbatim} plus the three \newenvironment lines, which the chapter preamble keeps
% verbatim from the project preamble). Nothing added, removed or redefined: these environments
% carry no mathematical content and the project itself typesets them as comments. Without these
% three lines a counted statement whose own text uses them fails with Environment slogan undefined.
\usepackage{verbatim}
\newenvironment{reference}{\comment}{\endcomment}
\newenvironment{slogan}{\comment}{\endcomment}
\newenvironment{history}{\comment}{\endcomment}
\usepackage{cleveref}
\crefname{originalproposition}{original proposition}{original propositions}
\Crefname{originalproposition}{Original proposition}{Original propositions}
\makeatletter
\newlabel{derived-proposition-left-derived-exists}{{0794}{source}{}{sourcepointer-1}{}}
\newlabel{derived-proposition-left-derived-exists@cref}{{[originalproposition][0][]0794}{[1][0][]source}{}{}{}}
\makeatother
\crefname{originalreference}{original reference}{original references}
\Crefname{originalreference}{Original reference}{Original references}
\makeatletter
\newlabel{lemma-O-homology-qis}{{08RX}{source}{}{sourcepointer-2}{}}
\newlabel{lemma-O-homology-qis@cref}{{[originalreference][0][]08RX}{[1][0][]source}{}{}{}}
\makeatother
\crefname{originalreference}{original reference}{original references}
\Crefname{originalreference}{Original reference}{Original references}
\makeatletter
\newlabel{lemma-compute-left-derived-pi-shriek}{{08PJ}{source}{}{sourcepointer-3}{}}
\newlabel{lemma-compute-left-derived-pi-shriek@cref}{{[originalreference][0][]08PJ}{[1][0][]source}{}{}{}}
\makeatother
\crefname{originalreference}{original reference}{original references}
\Crefname{originalreference}{Original reference}{Original references}
\makeatletter
\newlabel{lemma-compute-left-derived-pi-shriek-pre}{{08PI}{source}{}{sourcepointer-4}{}}
\newlabel{lemma-compute-left-derived-pi-shriek-pre@cref}{{[originalreference][0][]08PI}{[1][0][]source}{}{}{}}
\makeatother
\crefname{originallemma}{original lemma}{original lemmas}
\Crefname{originallemma}{Original lemma}{Original lemmas}
\makeatletter
\newlabel{sites-modules-lemma-module-quotient-flat}{{03EW}{source}{}{sourcepointer-5}{}}
\newlabel{sites-modules-lemma-module-quotient-flat@cref}{{[originallemma][0][]03EW}{[1][0][]source}{}{}{}}
\makeatother
\title{An isomorphism of the lower shrieks of the structure sheaves gives a base change formula for the lower shriek}
\author{The Stacks Project Authors}
\date{stacks-project:sites-cohomology.tex:local-snapshot-20261011; GNU FDL 1.2 or later}
\begin{document}
\maketitle
\noindent\textit{Editorial note.} All delivered mathematical text is the author's own, copied from the pinned Stacks Project chapter file sites-cohomology.tex; no statement or proof step was rewritten, reordered or removed.\par\medskip
\noindent\textit{Editorial note.} The counted proof is the author's complete argument: the proof is written out in full: the reduction to bounded above complexes with flat terms, the computation of both sides through the lower-shriek lemmas, and the final step on the fibre categories.\par\medskip
\noindent\textit{Editorial note.} The label lemma-base-change-by-qis names the counted statement, and the proof environment immediately following it in the source is delivered verbatim as its proof.\par\medskip
\noindent\textit{Editorial note.} The counted statement is carried by the author's own lemma environment in the chapter source.\par\medskip
\noindent\textit{Editorial note.} The statement and proof are copied from the pinned chapter file; the five cited inputs are kept as external pointers with their canonical Stacks tags.\par\medskip
\noindent\textit{Allowed standard prerequisites (editorial).}\begin{itemize}
\item Computation of the lower shriek through the cosimplicial resolution. Same chapter, the two lemmas computing the cohomology sheaves of L pi\_! for simplicial sheaves of rings. Supplies the procedure by which both sides of the identity are computed.
\item The colimit computation of the lower shriek and flat resolutions. Same chapter, the lemma computing L pi\_! of a module in terms of the fibre categories, together with Modules on Sites, the representation of a bounded above complex by one with flat terms. Supplies the reduction to a flat module and the final identification on the fibre categories.
\item The auxiliary computation for the restriction to a fibre category. Same chapter, the lemma comparing L pi\_! after extending scalars along a map of sheaves of rings. Supplies the result in the case of a single object of the simplicial category.
\end{itemize}\medskip
\ifdefined\hypertarget\hypertarget{sourcepointer-1}{}\fi
\noindent\textit{Source reference.} Original proposition 0794: sources/stacks-chow/tags.txt, source lines 9259--9259. Prerequisite pointer invoked at proof line(s) 10839; it names a standard fact of a sibling chapter that the counted argument uses, and every step of the counted proof built on it is retained verbatim in the delivered body.\par\medskip
\ifdefined\hypertarget\hypertarget{sourcepointer-2}{}\fi
\noindent\textit{Source reference.} Original reference 08RX: sources/stacks-sites-cohomology/pinned-original-source.tex, source lines 10422--10422. Prerequisite pointer invoked at proof line(s) 10852; it names a standard fact of the same chapter that the counted argument uses, and every step of the counted proof built on it is retained verbatim in the delivered body.\par\medskip
\ifdefined\hypertarget\hypertarget{sourcepointer-3}{}\fi
\noindent\textit{Source reference.} Original reference 08PJ: sources/stacks-sites-cohomology/pinned-original-source.tex, source lines 10615--10615. Prerequisite pointer invoked at proof line(s) 10846; it names a standard fact of the same chapter that the counted argument uses, and every step of the counted proof built on it is retained verbatim in the delivered body.\par\medskip
\ifdefined\hypertarget\hypertarget{sourcepointer-4}{}\fi
\noindent\textit{Source reference.} Original reference 08PI: sources/stacks-sites-cohomology/pinned-original-source.tex, source lines 10526--10526. Prerequisite pointer invoked at proof line(s) 10849; it names a standard fact of the same chapter that the counted argument uses, and every step of the counted proof built on it is retained verbatim in the delivered body.\par\medskip
\ifdefined\hypertarget\hypertarget{sourcepointer-5}{}\fi
\noindent\textit{Source reference.} Original lemma 03EW: sources/stacks-chow/tags.txt, source lines 4471--4471. Prerequisite pointer invoked at proof line(s) 10842; it names a standard fact of a sibling chapter that the counted argument uses, and every step of the counted proof built on it is retained verbatim in the delivered body.\par\medskip
\begin{lemma}
\label{lemma-base-change-by-qis}
Let $\mathcal{C}$ be a site. Let $\mathcal{A}_\bullet \to \mathcal{B}_\bullet$
be a homomorphism of simplicial sheaves of rings on $\mathcal{C}$.
If $L\pi_!\mathcal{A}_\bullet \to L\pi_!\mathcal{B}_\bullet$ is an
isomorphism in $D(\mathcal{C})$, then we have
$$
L\pi_!(K) =
L\pi_!(K \otimes^\mathbf{L}_{\mathcal{A}_\bullet} \mathcal{B}_\bullet)
$$
for all $K$ in $D(\mathcal{A}_\bullet)$.
\end{lemma}

\begin{proof}
Let $([n], U)$ be an object of $\Delta \times \mathcal{C}$. Since $L\pi_!$
commutes with colimits, it suffices to prove this for bounded above complexes
of $\mathcal{O}$-modules (compare with argument of
Derived Categories, Proposition \ref{derived-proposition-left-derived-exists}
or just stick to bounded above complexes). Every such complex is
quasi-isomorphic to a bounded above complex whose terms are flat modules, see
Modules on Sites, Lemma \ref{sites-modules-lemma-module-quotient-flat}.
Thus it suffices to prove the lemma for a flat $\mathcal{A}_\bullet$-module
$\mathcal{F}$. In this case the derived tensor product is the usual
tensor product and is a sheaf also. Hence by
Lemma \ref{lemma-compute-left-derived-pi-shriek}
we can compute the cohomology
sheaves of both sides of the equation by the procedure of
Lemma \ref{lemma-compute-left-derived-pi-shriek-pre}.
Thus it suffices to prove the result for the restriction of
$\mathcal{F}$ to the fibre categories (i.e., to $\Delta \times U$).
In this case the result follows from Lemma \ref{lemma-O-homology-qis}.
\end{proof}

\end{document}
